\documentclass[11pt,a4paper]{article}

\def\courseTitle{Industrial Security (Master)}
\def\sectionNumber{08}
\def\sectionTitle{Research Methods}
\def\sheetKind{Solutions}

% =====================================================================
% Shared preamble for exercise sheets and solution sheets.
%
% Define BEFORE \input{}-ing this file:
%   \def\courseTitle{...}
%   \def\sectionNumber{...}
%   \def\sectionTitle{...}
%   \def\sheetKind{Exercise Sheet}   or  Solution Sheet
%
% Optional:
%   \def\courseAuthor{...}
%
% The caller must issue \documentclass{...} (typically article) BEFORE
% loading this preamble.
% =====================================================================

\usepackage[utf8]{inputenc}
\usepackage[T1]{fontenc}
\usepackage[english]{babel}
\usepackage[a4paper, margin=2.2cm]{geometry}
\usepackage{graphicx}
\usepackage{listings}
\usepackage{xcolor}
\usepackage{tikz}
\usepackage{booktabs}
\usepackage{array}
\usepackage{enumitem}
\usepackage{titlesec}
\usepackage{fancyhdr}
\usepackage{hyperref}
\usepackage{amsmath}
\usepackage{amssymb}
\usepackage{tcolorbox}
\tcbuselibrary{skins, breakable}
\usepackage{needspace}

% Discourage solo lines split across pages.
\widowpenalty=10000
\clubpenalty=10000
\raggedbottom

% =====================================================================
% Shared TikZ styles for the Industrial Security lecture series.
% Loaded by every slide deck and exercise sheet that uses TikZ.
% =====================================================================

\usetikzlibrary{
  arrows.meta,
  shapes,
  shapes.geometric,
  shapes.symbols,
  shapes.multipart,
  positioning,
  fit,
  calc,
  backgrounds,
  decorations.pathmorphing,
  decorations.pathreplacing,
  decorations.text,
  patterns,
  matrix,
  shadows
}

% --- colour palette --------------------------------------------------
\definecolor{isOT}{HTML}{1F4E79}     % OT (deep blue)
\definecolor{isIT}{HTML}{6F2DA8}     % IT (purple)
\definecolor{isSafety}{HTML}{C00000} % safety red
\definecolor{isNet}{HTML}{2E7D32}    % network green
\definecolor{isAttack}{HTML}{B23A48} % attack red
\definecolor{isAccent}{HTML}{E08E0B} % amber
\definecolor{isMute}{HTML}{6B6B6B}   % grey
\definecolor{isLight}{HTML}{EDF2F8}  % light background
\definecolor{isPLC}{HTML}{2C5282}
\definecolor{isHMI}{HTML}{2F855A}
\definecolor{isSCADA}{HTML}{6B46C1}

% --- generic node styles --------------------------------------------
\tikzset{
  isnode/.style={
    draw, rounded corners=2pt, minimum height=8mm, minimum width=18mm,
    font=\small, align=center, thick
  },
  isbox/.style={isnode, fill=isLight},
  plc/.style={isnode, fill=isPLC!15, draw=isPLC, thick},
  hmi/.style={isnode, fill=isHMI!15, draw=isHMI, thick},
  scada/.style={isnode, fill=isSCADA!15, draw=isSCADA, thick},
  sensor/.style={isnode, fill=isNet!10, draw=isNet, minimum height=6mm, minimum width=14mm, font=\scriptsize},
  actuator/.style={isnode, fill=isAccent!15, draw=isAccent, minimum height=6mm, minimum width=14mm, font=\scriptsize},
  firewall/.style={isnode, fill=isAttack!15, draw=isAttack, thick},
  server/.style={isnode, fill=isMute!15, draw=isMute!80, thick},
  switch/.style={isnode, fill=isLight, draw=black!60, minimum height=6mm, minimum width=14mm, font=\scriptsize},
  workstation/.style={isnode, fill=white, draw=isOT, thick},
  attacker/.style={isnode, fill=isAttack!20, draw=isAttack, thick, font=\small\bfseries},
  inetnode/.style={draw, ellipse, minimum width=24mm, minimum height=10mm, font=\scriptsize, fill=isLight, thick},
  process/.style={isnode, fill=isOT!15, draw=isOT, thick, minimum width=24mm},
  zone/.style={
    draw, rounded corners=4pt, thick, dashed, inner sep=4mm
  },
  conduit/.style={-{Stealth[length=2.5mm]}, thick, isNet!80!black},
  attack/.style={-{Stealth[length=2.5mm]}, thick, isAttack, dashed},
  flow/.style={-{Stealth[length=2.5mm]}, thick, isOT!70!black},
  netline/.style={thick, black!60},
  axisline/.style={-{Stealth[length=2mm]}, thick, black!70},
  tag/.style={font=\scriptsize\itshape, fill=white, inner sep=1pt},
  level/.style={
    draw, fill=isLight, minimum width=0.85\linewidth, minimum height=8mm,
    font=\small, anchor=west
  },
  brace/.style={decorate, decoration={brace, amplitude=4pt, raise=2pt}},
  legendbox/.style={draw, rounded corners=2pt, fill=isLight, inner sep=2mm, font=\scriptsize, align=left}
}

% --- handy macros ----------------------------------------------------
% \plcicon, \hmiicon etc as simple labelled boxes; useful inline.
\newcommand{\plcicon}[1][PLC]{\tikz[baseline=-0.6ex]{\node[plc, minimum width=10mm, minimum height=5mm, font=\scriptsize, inner sep=1pt]{#1};}}
\newcommand{\hmiicon}[1][HMI]{\tikz[baseline=-0.6ex]{\node[hmi, minimum width=10mm, minimum height=5mm, font=\scriptsize, inner sep=1pt]{#1};}}
\newcommand{\fwicon}[1][FW]{\tikz[baseline=-0.6ex]{\node[firewall, minimum width=10mm, minimum height=5mm, font=\scriptsize, inner sep=1pt]{#1};}}


\providecommand{\courseAuthor}{Industrial Security Lecture Series}
\providecommand{\courseInstitute}{Department of Computer Science}
\providecommand{\companyLogo}{logo}

% --- Corporate branding overrides ------------------------------------
% Picks up logo / author / brand colours from the repo-root branding.tex
% when present; falls back to the defaults above otherwise.
\IfFileExists{../../../branding.tex}{% =====================================================================
% Corporate / institutional branding -- single file to customise the
% look of the entire course (slides, notes, exercises, labs).
%
% HOW TO REBRAND IN UNDER A MINUTE:
%   1. Replace `branding/logo.pdf` with your own logo
%      (PDF preferred; PNG and JPG also work -- name it `logo.<ext>`).
%   2. (Optional) change the two colours below to your brand palette.
%   3. (Optional) override the author/institute lines below.
%   4. Re-run `make` at the repo root. That's it.
%
% Everything in this file has sensible defaults; you can delete a line
% to fall back to the built-in defaults, or comment it out.
%
% This file is loaded automatically by the slides, exercise, and lab
% preambles -- you never need to \input it manually.
% =====================================================================

% --- Logo --------------------------------------------------------------
% Path is resolved from each leaf-build directory; the default points at
% the shared `branding/` folder at the repo root. If you put your logo
% somewhere else, set the full or relative path here (without extension):
\renewcommand{\companyLogo}{../../../branding/logo}

% --- Author / institute (appears on the title page footer) ------------
\renewcommand{\courseAuthor}{}
\renewcommand{\courseInstitute}{}

% --- Brand colours ----------------------------------------------------
% slideAccent      = primary brand colour (title bands, accent rule)
% slideSecondary   = secondary brand colour (section badge, highlight)
% Both use HTML hex codes. Keep enough contrast against white text.
\definecolor{slideAccent}{HTML}{00205B}      % deep navy
\definecolor{slideSecondary}{HTML}{00A0DC}   % light blue accent
}{}

% --- listings -------------------------------------------------------
\definecolor{codebg}{rgb}{0.96,0.96,0.96}
\lstset{
  backgroundcolor=\color{codebg},
  basicstyle=\ttfamily\footnotesize,
  breaklines=true,
  frame=single,
  numbers=left,
  numberstyle=\tiny\color{black!50},
  showstringspaces=false,
  tabsize=2
}

% --- header / footer ------------------------------------------------
\pagestyle{fancy}
\fancyhf{}
\renewcommand{\headrulewidth}{0.4pt}
\renewcommand{\footrulewidth}{0pt}
\lhead{\small \courseTitle}
\rhead{\small Section \sectionNumber{} -- \sheetKind}
\cfoot{\thepage}

% --- task / solution boxes -----------------------------------------
% `before={...}` guards every task/solution box with \needspace so the
% box doesn't start with only one or two lines of breathing room at the
% bottom of a page. If less than ~9 lines remain, LaTeX bumps the box
% to the next page; otherwise it starts here and breaks normally if it
% runs long.
% Boxes are `breakable` so very long solutions don't blow the page,
% but `before={\needspace{14\baselineskip}}` pushes them to a fresh
% page whenever fewer than ~14 lines remain. That covers the vast
% majority of single-question splits in practice.
\newtcolorbox{taskbox}[1]{
  enhanced, breakable,
  colback=isLight, colframe=isOT,
  fonttitle=\bfseries\color{white},
  coltitle=white,
  colbacktitle=isOT,
  title={#1},
  arc=2pt, boxrule=0.6pt,
  left=2mm, right=2mm, top=2mm, bottom=2mm,
  before={\par\vspace{2mm}\needspace{14\baselineskip}\par},
  after={\par\vspace{2mm}}
}

\newtcolorbox{solutionbox}{
  enhanced, breakable,
  colback=isNet!5, colframe=isNet,
  fonttitle=\bfseries\color{white},
  coltitle=white,
  colbacktitle=isNet,
  title={Solution},
  arc=2pt, boxrule=0.6pt,
  left=2mm, right=2mm, top=2mm, bottom=2mm,
  before={\par\vspace{2mm}\needspace{14\baselineskip}\par},
  after={\par\vspace{2mm}}
}

\newtcolorbox{hintbox}{
  enhanced, breakable,
  colback=isAccent!10, colframe=isAccent,
  fonttitle=\bfseries\color{white}, coltitle=white, colbacktitle=isAccent,
  title={Hint},
  arc=2pt, boxrule=0.6pt
}

\newtcolorbox{warnbox}{
  enhanced, breakable,
  colback=isAttack!8, colframe=isAttack,
  fonttitle=\bfseries\color{white}, coltitle=white, colbacktitle=isAttack,
  title={Caution},
  arc=2pt, boxrule=0.6pt
}

% --- section formatting --------------------------------------------
\titleformat{\section}{\Large\bfseries\color{isOT}}{\thesection}{0.5em}{}
\titleformat{\subsection}{\large\bfseries\color{isOT!80!black}}{\thesubsection}{0.5em}{}

% --- title block macro ----------------------------------------------
\newcommand{\sheetheader}{%
  \begin{center}
    {\Large\bfseries \courseTitle}\\[0.3em]
    {\large \sheetKind{} -- Section \sectionNumber}\\[0.2em]
    {\large \sectionTitle}\\[0.2em]
    \rule{\linewidth}{0.4pt}
  \end{center}
  \vspace{0.5em}
}


\begin{document}
\sheetheader

\noindent
\textit{How to use this sheet:} the solutions are model answers, not the only
correct ones. Each box ends with a \textbf{Rubric} listing what must appear for
full marks; a different but defensible answer that covers the rubric is correct.
Research-methods marking rewards a \emph{narrow, falsifiable, honestly-reported}
answer over an ambitious vague one.

\begin{solutionbox}
\textbf{Exercise 8.1 -- Vague topic to falsifiable RQ.}
\textbf{RQ:} ``Does the Modbus/TCP write-coil injection rate (writes/s) change the
per-minute count of unsolicited-write notices raised by a Zeek monitor on the
Bachelor OpenPLC stack?'' \textbf{H1:} a higher injection rate increases the
detection count by a measurable margin. \textbf{H0:} injection rate has no
measurable effect on the detection count. \textbf{Scope IN:} one protocol
(Modbus/TCP), one component (OpenPLC + Zeek), one metric (notices/min). \textbf{Scope OUT:}
all other protocols, real plants, evasion variants.
\textit{FINER:} \emph{Feasible} --- runs in one session on hardware you own;
\emph{Interesting} --- detection scaling informs sensor tuning; \emph{Novel} ---
done on a fully reproducible, pinned testbed rather than a private rig;
\emph{Ethical} --- testbed-only, no third party; \emph{Relevant} --- write-coil
injection is a real OT manipulation class.
\textbf{Rubric:} exactly one RQ, one H1, one H0; H0 is the genuine null (``no
effect''), not a strawman; scope has both an IN and an OUT line; all five FINER
criteria addressed and tied to \emph{this} RQ, not generic.
\end{solutionbox}

\begin{solutionbox}
\textbf{Exercise 8.2 -- Threats the abstract misses.}
\textit{Internal:} train and test traffic came from the same one-afternoon
capture, so the model may have memorised session artefacts (IP/port/timing) rather
than attack behaviour --- no train/test split by time or run is described.
\textit{External:} a single testbed (SWaT), one protocol window, one afternoon
cannot support ``ready for deployment in real plants''; generalisation to other
vendors/sites is asserted, not shown. \textit{Construct:} ``accuracy'' does not
measure ``detects injection'' on a class-imbalanced stream --- if 99.7\,\% of
packets are benign, a model that flags nothing scores $\approx$99.7\,\% accuracy.
\textit{Statistical/conclusion red flag:} the headline is a single accuracy number
with no precision/recall/F1, no confusion matrix, no class balance, no variance
over runs. \textit{Should have reported:} time-disjoint train/test split;
precision, recall, F1 (or PR-AUC) per class; the base rate / class balance; a
cross-testbed or held-out-attack evaluation; spread over repeated runs; and a
claim scoped to the testbed, not to ``real plants''.
\textbf{Rubric:} one each of internal, external, construct named correctly; the
imbalanced-accuracy red flag identified; at least one concrete ``report instead''
fix; ``ready for deployment'' flagged as unsupported.
\end{solutionbox}

\begin{hintbox}
For 8.2 and 8.6, the reflex question is always \emph{``compared to what, measured
how, with what spread, and would it survive a different day/testbed?''} Accuracy on
imbalanced data and a difference-of-means with no variance are the two classic OT-paper traps.
\end{hintbox}

\begin{solutionbox}
\textbf{Exercise 8.3 -- Reproducible experiment design.}
\vspace{0.5mm}
{\footnotesize
\par\begin{tabular}{@{}p{4.0cm}p{11.4cm}@{}}
\toprule
\textbf{Element} & \textbf{Choice} \\
\midrule
Independent variable & Modbus write-coil rate, levels \{1, 10, 50\} writes/s \\
Dependent variable & Zeek unsolicited-write notices per minute (count/min) \\
Controlled variables & stack commit hash, container CPU/mem limits, payload size, MTU, target coil range, capture interface \\
Baseline / control & idle poll loop only, zero injection \\
Metric definition & notices/min = \texttt{grep -c} of the notice tag over a fixed 60\,s window \\
Repetitions \& order & $\geq 5$ runs per level; level order randomised to defeat warm-up bias \\
Confound 1 + fix & Docker noisy-neighbour CPU contention $\to$ pin CPU set, no other load \\
Confound 2 + fix & host clock drift across runs $\to$ NTP-sync host; timestamp every command \\
\bottomrule
\end{tabular}}
\vspace{0.5mm}
\textit{Reject H0 if} the mean notices/min rises monotonically and significantly
with rate across levels (e.g.\ a trend test, or non-overlapping CIs between 1 and
50 writes/s); otherwise H0 stands and ``no measurable effect'' is reported as the
result. \textbf{Rubric:} every table row present; a real baseline (idle, not just
``low rate''); $\geq 5$ reps and randomised order; at least two confounds each with
a concrete mitigation; an explicit decision rule for rejecting H0.
\end{solutionbox}

\begin{solutionbox}
\textbf{Exercise 8.4 -- Citation matrix + classification.}
\vspace{0.5mm}
{\footnotesize
\par\begin{tabular}{@{}p{3.1cm}p{2.4cm}p{2.6cm}p{1.8cm}p{3.8cm}@{}}
\toprule
\textbf{Paper} & \textbf{Method} & \textbf{Substrate} & \textbf{Artefact?} & \textbf{Gap vs.\ my RQ} \\
\midrule
Shodan census '23 & measurement & Internet scan & dataset maybe & no detection, no testbed, no causal claim \\
TRITON '18 & case study & one real plant & no & single incident; no reproducible metric \\
GooseHunter '22 & design science & testbed/sim & tool likely & GOOSE not Modbus; reproducibility unstated \\
Scan-cycle jitter '24 & experiment & controlled testbed & code likely & jitter side channel, not write-injection notices \\
TTP taxonomy '21 & theory building & campaign reports & no & descriptive; no measurement on a stack \\
\bottomrule
\end{tabular}}
\vspace{0.5mm}
\textit{Where my contribution sits:} a controlled, fully pinned and re-runnable
study of Modbus write-injection \emph{detection scaling} on an open testbed --- the
empty cell across all five rows.
\textbf{Rubric:} method correctly classified for each (census = measurement, TRITON
= case study, GooseHunter = design science, jitter = experiment, taxonomy = theory
building); substrate and artefact columns sensible; the gap column is specific to
the RQ, not generic; the contribution statement points at the empty cell.
\end{solutionbox}

\begin{solutionbox}
\textbf{Exercise 8.5 -- Ethics / disclosure / dual-use statement.}
\textit{Scope of harm:} the crash was reproduced only on a spare unit the
researcher owns, on an isolated bench network; no third-party system, no
production plant, no Internet-reachable target was touched. \textit{Beneficence vs.\
maleficence (Menlo):} asset owners and the vendor benefit from a fix before the
flaw is found independently and weaponised; the harm of disclosure is bounded by
withholding the working crasher until a patch exists; the balance is positive
because the same family is widely deployed and a denial-of-service on a controller
can endanger a physical process. \textit{Coordinated disclosure:} notify the vendor
PSIRT first with a reproducible report and a proposed timeline, and the
coordinating CERT (CISA / BSI / JPCERT, per region) in parallel for a CVE and
tracking. Use an ICS-realistic clock --- ~180~days with a 60-day grace --- rather
than Project Zero's 90/30, justified by ICS maintenance windows and
recertification; if the vendor goes silent or refuses, escalate through the CERT
and disclose on the published deadline. \textit{Dual-use:} publish the trigger
condition class, the detection signature, and the mitigation (input validation /
filtering on function-code 90); withhold the turnkey packet and any auto-exploit
until patch availability. \textit{Standards relied on:} ISO/IEC 29147:2018
(disclosure), ISO/IEC 30111:2019 (handling), the Menlo Report, CERT/CC policy.
\textbf{Rubric:} testbed-only scope stated; explicit benefit/harm argument with
Menlo framing; named recipients (vendor PSIRT + coordinating CERT); an
ICS-appropriate timeline with deviation justified; clear publish-vs-withhold split;
at least two real standards cited.
\end{solutionbox}

\begin{solutionbox}
\textbf{Exercise 8.6 -- Results-claim soundness.}
The claim is unsupported as written. (1)~\emph{Sample size:} $n=5$ is small; with
high run-to-run variance the means may be statistically indistinguishable.
(2)~\emph{No spread reported:} 0.91 vs.\ 0.88 means nothing without SD or a
confidence interval --- if SDs are $\pm 0.04$ the intervals overlap heavily.
(3)~\emph{No test:} a difference of means must be tested (paired
$t$-test or a non-parametric Wilcoxon on the paired runs), with the $p$-value and
the test named. (4)~\emph{Effect size vs.\ significance:} a 0.03 F1 gap may be
statistically real yet operationally trivial; report the effect size (e.g.\
Cohen's $d$) and ask whether it matters in deployment. (5)~\emph{Cherry-picking:}
F1 alone may hide a precision/recall trade-off; check the metric was pre-chosen, not
selected after seeing results (HARKing), and correct for multiple comparisons if
several metrics were tried. (6)~``Proves'' is wrong --- evidence \emph{supports}
or \emph{fails to reject}; nothing is proven. \textit{Would require before
believing:} per-run paired data, mean $\pm$ CI, a named significance test with its
$p$-value, an effect size, and the same evaluation protocol for both systems.
\textbf{Rubric:} at least four of the six points; recognises overlapping
variance / no test as the core flaw; distinguishes significance from effect size;
rejects ``proves''; states a concrete bar for belief.
\end{solutionbox}

\begin{solutionbox}
\textbf{Exercise 8.7 -- Reproducibility package outline.}
\textit{Layout:} \texttt{README.md} (claim + cold-start command block);
\texttt{scripts/} (data collection + \texttt{plot\_results.py}, no manual edits);
\texttt{data/} (raw CSV or fetch script, with provenance); \texttt{results/} +
\texttt{figures/} (regenerated outputs + \texttt{CHECKSUMS}); \texttt{requirements.txt}
(\texttt{pip freeze}, exact versions); \texttt{docs/env.txt} (Docker + compose
versions, pinned Bachelor-stack commit hash); \texttt{LICENSE} (permissive).
\textit{Three ICS reproducibility killers + the pin that defeats each:}
(1)~unpinned dependencies $\to$ \texttt{requirements.txt} from \texttt{pip freeze}
and a fixed base image digest; (2)~``works only on the author's testbed'' $\to$
ship against the public Bachelor Docker stack pinned to a recorded commit hash;
(3)~moving/closed data $\to$ commit the raw CSV (or a hashed fetch script) with a
\texttt{CHECKSUMS} file so a re-runner confirms a byte match. \textit{One-command
path:} \texttt{git clone \dots\ \&\& docker compose up -d \&\& ./scripts/run\_all.sh}
$\to$ \texttt{figures/result.png}, hash-checked against \texttt{CHECKSUMS}.
\textbf{Rubric:} maps to ACM Available + Reusable (env pinning + runnable, not just
posted); names the three killers from the lecture (closed deps, single-site
equipment, moving data) each with a concrete pin; gives a genuine clone-to-figure
command path with a checksum verification.
\end{solutionbox}

\begin{solutionbox}
\textbf{Exercise 8.9 -- Method selection.}
\begin{enumerate}\setlength{\itemsep}{1pt}\small
  \item IEC~104 endpoint count --- \textbf{measurement}: the question is ``how
        many'', answered by an Internet/Shodan census. \emph{Threat:} coverage /
        sampling bias (scanners miss firewalled or rate-limited hosts), so the
        count is a lower bound.
  \item Why Maroochy succeeded --- \textbf{case study}: a single rich incident, ``what
        happened and why'', from forensic and organisational record. \emph{Threat:}
        external validity --- one disgruntled-insider case does not generalise to
        all OT.
  \item GOOSE + MAC latency --- \textbf{experiment}: ``does X cause Y'' with a
        controlled IV (MAC on/off) and a measured DV (latency). \emph{Threat:}
        construct/internal --- lab latency may not equal substation-bus latency, and
        confounds (CPU load) must be controlled.
  \item S7comm fuzzer --- \textbf{design science}: build an artefact (the fuzzer)
        and evaluate it by bugs found. \emph{Threat:} the evaluation is only as
        good as the corpus/oracle --- ``found no bugs'' may mean a weak fuzzer, not
        a safe stack.
\end{enumerate}
\textbf{Rubric:} each method correct (measurement / case study / experiment /
design science); two-sentence justification ties method to the question form; the
named threat is the one that method specifically invites.
\end{solutionbox}

\begin{solutionbox}
\textbf{Exercise 8.8 -- Disclosure-timeline essay (grading notes).}
A strong essay takes a clear position and defends it with the named trade-off, not
slogans. It must acknowledge that ICS patching genuinely differs from IT ---
scheduled maintenance windows, safety recertification, heterogeneous fleets, and
legacy assets that may be \emph{unpatchable} --- which pushes toward a longer clock
(e.g.\ 180/60) and a coordinating CERT (CISA / BSI / JPCERT) under ISO/IEC
29147/30111. But it must also concede the cost of an over-long embargo: users sit
exposed, an independent finder may already hold the bug, and a slow vendor loses all
pressure to act. The best answers resolve the tension with a \emph{conditional}
policy: a default ICS window longer than 90/30, mitigations published early even
when the full patch is slow, a hard deadline regardless of vendor cooperation, and
an explicit escalation when the vendor goes silent (CERT escalation, then disclose
on the published date). \textbf{Rubric:} a defensible thesis stated; engages ICS
patch-rollout reality \emph{and} the over-embargo cost (both sides); names a
coordinating CERT and at least two standards/policies (ISO/IEC 29147, 30111,
CERT/CC, Project Zero); states a concrete plan for vendor silence; 250--350 words.
Deduct for a one-sided ``always wait longer'' with no deadline.
\end{solutionbox}

\end{document}
